\begin{workedexample}[Worked Example: Class efficiency limits]
Ideal drain efficiency depends on conduction angle. Class~A conducts the
whole cycle ($360^{\circ}$) with a theoretical maximum of $50\%$; Class~B
conducts a half cycle ($180^{\circ}$) with $\pi/4 = 78.5\%$; Class~AB sits
between. Switching classes (D, E, F) can approach $100\%$ ideally but require
hard switching. Backing a linear PA off $6\ \text{dB}$ from saturation trades
much of that efficiency for low distortion.
\end{workedexample}

\begin{keytakeaways}
\begin{itemize}[leftmargin=1.4em,itemsep=2pt]
\item Efficiency rises as conduction angle falls: A ($50\%$) $<$ AB $<$ B ($78.5\%$) $<$ switching.
\item Linearity and efficiency trade off; back-off buys linearity at lower efficiency.
\item Doherty and digital predistortion recover efficiency/linearity for modulated signals.
\end{itemize}
\end{keytakeaways}

\begin{exercises}
\item What is the ideal maximum efficiency of Class A?
\item What is the ideal maximum efficiency of Class B?
\item Backing a PA off from saturation improves what, at what cost?
\end{exercises}

\furtherreading{Cripps~\cite{cripps}; Pozar~\cite{pozar} (ch.~12).}
